\documentclass[a4paper,11pt]{article}

\usepackage[utf8]{inputenc}
\usepackage[T1]{fontenc}
\usepackage[ngerman]{babel}
\usepackage{amsmath,amssymb}
\usepackage{geometry}
\geometry{margin=2.0cm}
\usepackage{enumitem}
\usepackage{fancyhdr}
\usepackage{xcolor}
\usepackage{tikz}
\usepackage{titlesec}
\usepackage{array}
\usepackage{colortbl}

%================= Dark-Mode Farben (Catppuccin Mocha) =================
\definecolor{base}{HTML}{1E1E2E}     % Hintergrund
\definecolor{fg}{HTML}{CDD6F4}       % Text (heller Vordergrund)
\definecolor{accent}{HTML}{89B4FA}   % Überschriften / Akzent
\definecolor{accent2}{HTML}{F5C2E7}  % Punkteangabe
\definecolor{gridcolor}{HTML}{45475A}% Karo-Gitter / Linien (dezent)

\pagecolor{base}

\newcommand{\transp}{^{\mathsf{T}}}
\newcommand{\norm}[1]{\left\lVert #1 \right\rVert}
\newcommand{\rang}{\operatorname{rang}}
\newcommand{\diag}{\operatorname{diag}}

% Punkteangabe rechtsbündig an der Aufgabenüberschrift
\newcommand{\pkt}[1]{\hfill\textnormal{\color{accent2}\itshape (#1)}}

% Karierte Antwortfläche (Karopapier) -- #1 = Höhe des Bereichs (z. B. 5cm).
% Durchgehendes Karo-Raster, keine einzelnen Schreiblinien, kein \hrulefill.
\newcommand{\karo}[1]{\par\vspace{0.3cm}\noindent
  \begin{tikzpicture}
    \draw[step=5mm, gridcolor, very thin] (0,0) grid (\linewidth,#1);
  \end{tikzpicture}\par\vspace{0.3cm}}

% Ankreuzkästchen
\newcommand{\kasten}{{\color{gridcolor}$\square$}}

%================= Überschriften in Akzentfarbe =================
\titleformat{\section}{\color{accent}\normalfont\Large\bfseries}{}{0pt}{}

%================= Kopf-/Fußzeile =================
\pagestyle{fancy}
\fancyhf{}
\lhead{\color{fg}\small FLA \& Analytische Geometrie}
\rhead{\color{fg}\small Probeklausur 01}
\cfoot{\color{fg}\thepage}
\renewcommand{\headrule}{{\color{gridcolor}\hrule height 0.4pt}}

\arrayrulecolor{gridcolor}
\setlength{\parindent}{0pt}

\begin{document}
\color{fg}

\thispagestyle{fancy}

\begin{center}
  {\color{accent}\LARGE\bfseries Probeklausur}\\[0.3em]
  {\color{fg}\large Fortgeschrittene Lineare Algebra und Analytische Geometrie}\\[0.6em]
  {\color{gridcolor}\rule{\linewidth}{0.4pt}}
\end{center}

\vspace{0.4em}

% ---------------- Kopf: Name / Datum / Rahmendaten ----------------
\renewcommand{\arraystretch}{1.6}
\begin{tabular}{@{}p{0.5\linewidth}@{}p{0.5\linewidth}@{}}
  Name: \ \rule[-2pt]{0.7\linewidth}{0.4pt} &
  Datum: \ \rule[-2pt]{0.55\linewidth}{0.4pt}
\end{tabular}

\vspace{0.6em}

\begin{tabular}{@{}l p{0.72\linewidth}@{}}
  \textbf{\color{accent}Bearbeitungszeit:} & 90 Minuten \\
  \textbf{\color{accent}Gesamtpunkte:}     & 60 Punkte \\
  \textbf{\color{accent}Hilfsmittel:}      & ein beidseitig handbeschriebenes Formelblatt (DIN A4), nicht-programmierbarer Taschenrechner \\
\end{tabular}

\vspace{0.4em}

\textit{\color{fg}\small Hinweis: Alle Ergebnisse sind durch nachvollziehbare Rechenwege zu begründen. Notieren Sie Zwischenschritte.}

\vspace{1.0em}

% ---------------- Punktetabelle ----------------
\renewcommand{\arraystretch}{1.5}
\noindent
\begin{tabular}{|c|*{9}{c|}c|}
\hline
\textbf{Aufgabe} & 1 & 2 & 3 & 4 & 5 & 6 & 7 & 8 & 9 & \textbf{$\Sigma$} \\
\hline
\textbf{max.\ Punkte} & 5 & 6 & 9 & 6 & 8 & 7 & 6 & 5 & 8 & \textbf{60} \\
\hline
\textbf{erreicht} & & & & & & & & & & \\
\hline
\end{tabular}

\renewcommand{\arraystretch}{1.0}

% ================================================================
\clearpage
\section*{Aufgabe 1 -- Quickies \pkt{5 Punkte}}

Sind die folgenden Aussagen \emph{wahr} oder \emph{falsch}? Kreuzen Sie jeweils an.

\vspace{0.6em}
\begin{enumerate}[label=\alph*)]
  \item Jede reelle symmetrische Matrix ist diagonalisierbar.
        \hfill \kasten\ wahr \quad \kasten\ falsch
  \item Für eine orthogonale Matrix $Q$ gilt $Q^{-1} = Q\transp$.
        \hfill \kasten\ wahr \quad \kasten\ falsch
  \item Jede invertierbare Matrix ist positiv definit.
        \hfill \kasten\ wahr \quad \kasten\ falsch
  \item Die Determinante einer oberen Dreiecksmatrix ist das Produkt ihrer Diagonaleinträge.
        \hfill \kasten\ wahr \quad \kasten\ falsch
  \item Die Singulärwerte einer Matrix können negativ sein.
        \hfill \kasten\ wahr \quad \kasten\ falsch
\end{enumerate}

\karo{5cm}

% ================================================================
\clearpage
\section*{Aufgabe 2 -- CR-Zerlegung \pkt{6 Punkte}}

Gegeben sei
\[
A = \begin{pmatrix} 1 & 2 & 4 \\ 1 & 3 & 5 \\ 2 & 5 & 9 \end{pmatrix}.
\]
\begin{enumerate}[label=\alph*)]
  \item Bestimmen Sie den Rang von $A$. \pkt{2}
  \item Wählen Sie eine Basis des Spaltenraums. \pkt{1}
  \item Bestimmen Sie die CR-Zerlegung $A = CR$. \pkt{3}
\end{enumerate}

\karo{12cm}

% ================================================================
\clearpage
\section*{Aufgabe 3 -- LU-Zerlegung \pkt{9 Punkte}}

Gegeben sei
\[
A = \begin{pmatrix} 2 & 1 & 1 \\ 4 & 3 & 5 \\ 2 & 4 & 12 \end{pmatrix}.
\]
\begin{enumerate}[label=\alph*)]
  \item Bestimmen Sie eine LU-Zerlegung. \pkt{4}
  \item Bestimmen Sie $\det(A)$. \pkt{1}
  \item Berechnen Sie $A^{-1}$. \pkt{4}
\end{enumerate}

\karo{14cm}

% ================================================================
\clearpage
\section*{Aufgabe 4 -- Cholesky-Zerlegung \pkt{6 Punkte}}

Gegeben sei
\[
A = \begin{pmatrix} 9 & 3 & 6 \\ 3 & 5 & 4 \\ 6 & 4 & 14 \end{pmatrix}.
\]
\begin{enumerate}[label=\alph*)]
  \item Prüfen Sie die positive Definitheit. \pkt{3}
  \item Bestimmen Sie die Cholesky-Zerlegung $A = L L\transp$. \pkt{3}
\end{enumerate}

\karo{13cm}

% ================================================================
\clearpage
\section*{Aufgabe 5 -- QR-Zerlegung \pkt{8 Punkte}}

Gegeben seien
\[
Q = \frac{1}{5}\begin{pmatrix} 3 & -4 \\ 4 & 3 \end{pmatrix},
\qquad
R = \begin{pmatrix} 5 & 10 \\ 0 & 5 \end{pmatrix}.
\]
Weiter sei $A = QR$ und
\[
b = \begin{pmatrix} 7 \\ 26 \end{pmatrix}.
\]
\begin{enumerate}[label=\alph*)]
  \item Bestimmen Sie $A = QR$. \pkt{2}
  \item Lösen Sie das Gleichungssystem $Ax = b$ mithilfe der QR-Zerlegung. \pkt{4}
  \item Bestimmen Sie $Q^{-1}$. \pkt{2}
\end{enumerate}

\karo{13cm}

% ================================================================
\clearpage
\section*{Aufgabe 6 -- Eigenwerte und Diagonalisierung \pkt{7 Punkte}}

Gegeben sei
\[
A = \begin{pmatrix} 0 & 1 \\ 2 & 1 \end{pmatrix}.
\]
\begin{enumerate}[label=\alph*)]
  \item Bestimmen Sie die Eigenwerte. \pkt{3}
  \item Bestimmen Sie Eigenvektoren. \pkt{2}
  \item Geben Sie $V$ und $D$ mit $A = V D V^{-1}$ an. \pkt{2}
\end{enumerate}

\karo{13cm}

% ================================================================
\clearpage
\section*{Aufgabe 7 -- SVD und Pseudoinverse \pkt{6 Punkte}}

Gegeben sei
\[
A = \begin{pmatrix} 4 & 0 \\ 0 & 0 \\ 0 & 3 \end{pmatrix}.
\]
\begin{enumerate}[label=\alph*)]
  \item Berechnen Sie $A\transp A$. \pkt{2}
  \item Bestimmen Sie die Singulärwerte. \pkt{2}
  \item Bestimmen Sie die Pseudoinverse $A^{+}$. \pkt{2}
\end{enumerate}

\karo{12cm}

% ================================================================
\clearpage
\section*{Aufgabe 8 -- Hyperebene \pkt{5 Punkte}}

Gegeben sei die Hyperebene
\[
2x_1 - x_2 + 2x_3 = 4.
\]
\begin{enumerate}[label=\alph*)]
  \item Bestimmen Sie einen Normalenvektor. \pkt{1}
  \item Liegt $p = (1,2,3)\transp$ auf der Ebene? \pkt{2}
  \item Berechnen Sie den Abstand von $p$ zur Ebene. \pkt{2}
\end{enumerate}

\karo{12cm}

% ================================================================
\clearpage
\section*{Aufgabe 9 -- Finite Differenzen \pkt{8 Punkte}}

Gegeben sei das Randwertproblem
\[
-u''(x) = 6, \qquad u(0) = u(1) = 0.
\]
Verwenden Sie die inneren Gitterpunkte
\[
x_1 = 0{,}25, \quad x_2 = 0{,}5, \quad x_3 = 0{,}75.
\]
\begin{enumerate}[label=\alph*)]
  \item Leiten Sie mit dem zentralen Differenzenquotienten aus der Differential\-gleichung das lineare Gleichungssystem $Ax = b$ her. \pkt{6}
  \item Nennen Sie vier strukturelle Eigenschaften der Matrix $A$. \pkt{2}
\end{enumerate}

\karo{13cm}

\end{document}
